\documentclass[12pt,letterpaper]{report}

\usepackage[utf8]{inputenc}
\usepackage[T1]{fontenc}
\usepackage[spanish,es-nodecimaldot]{babel}
\usepackage[letterpaper,margin=2.5cm]{geometry}
\usepackage{array}
\usepackage{booktabs}
\usepackage{tabularx}
\usepackage{xcolor}
\usepackage{hyperref}
\usepackage{enumitem}

\hypersetup{
  colorlinks=true,
  linkcolor=blue!50!black,
  urlcolor=blue!60!black,
  pdftitle={Análisis de arquitectura y post-mortem grupal -- Bosconian},
  pdfauthor={Equipo Bosconian Recreation}
}

\newcommand{\archivo}[1]{\texttt{\detokenize{#1}}}
\renewcommand{\arraystretch}{1.2}

\begin{document}
\hypersetup{pageanchor=false}
\begin{titlepage}
\centering
\vspace*{1.2cm}
{\Huge\bfseries Análisis de arquitectura\\[0.35em] y post-mortem grupal\par}
\vspace{1.0cm}
{\LARGE Bosconian Recreation\par}
\vspace{0.25cm}
{\large Recreación de Bosconian para Sharp X68000\par}
\vspace{0.5cm}
{\large Introducción al Desarrollo de Videojuegos\par}
\vfill
\begin{tabular}{>{\centering\arraybackslash}p{0.42\textwidth} >{\centering\arraybackslash}p{0.42\textwidth}}
\textbf{Integrante 1} & \textbf{Integrante 2} \\
Manuel Santiago Prieto Hernández & Daniela Munar Gonzalez \\
[1.4em]
\textbf{Integrante 3} & \textbf{Integrante 4} \\
Camilo Puerto Peña & Camilo Inoue Ruda \\
\end{tabular}
\vfill
{\large 23 de noviembre de 2026\par}
\end{titlepage}
\hypersetup{pageanchor=true}
\pagenumbering{roman}
\tableofcontents
\clearpage
\pagenumbering{arabic}

\chapter*{Alcance}
\addcontentsline{toc}{chapter}{Alcance}
Este documento presenta el análisis de la arquitectura y el post-mortem del proyecto Bosconian Recreation. La descripción técnica se basa en el enunciado, la documentación y el código del repositorio.

\chapter{Resumen del proyecto}

El proyecto consiste en desarrollar una recreación jugable de \emph{Bosconian}, tomando como referencia la versión de Sharp X68000. Se eligió Python con Pygame y Esper para trabajar con una ventana de 256\,$\times$\,208 píxeles y una arquitectura de entidades, componentes y sistemas (ECS).

La arquitectura separa el motor y las escenas de los datos de las entidades y de los sistemas que actualizan el juego. Los niveles y diversos parámetros se leen desde archivos JSON. En la versión examinada existen el jugador, enemigos, bases, torretas, obstáculos, colisiones, cámara, minimapa, escenas de menú y diez configuraciones de nivel. El proyecto incluye, además, recursos gráficos y de audio.

La implementación cuenta con la base técnica y un conjunto amplio de sistemas jugables. En la versión de \archivo{main} examinada, aún deben completarse requisitos obligatorios: activar el avance de ronda al destruir todas las bases, implementar las alertas amarilla y roja, mostrar las bases en el minimapa, persistir el récord y otorgar vidas extra según el puntaje. También debe ajustarse la habilidad especial para que corresponda con la especificación de cuatro proyectiles.

\chapter{Evolución del desarrollo}

El repositorio registra una evolución incremental: primero se establecieron el motor y la base ECS; posteriormente se incorporaron la compatibilidad con Esper 2, los recursos, la cámara, los controles del jugador, el combate y las configuraciones de niveles. La tabla relaciona el plan de seis semanas con los avances visibles en el código. Dado que el plan utiliza semanas relativas y el historial no asigna cada cambio a una semana específica, las etapas representan la secuencia propuesta y no fechas históricas exactas.

\begin{center}\small
\begin{tabularx}{\textwidth}{@{}>{\raggedright\arraybackslash}p{0.12\textwidth} >{\raggedright\arraybackslash}p{0.31\textwidth} >{\raggedright\arraybackslash}X@{}}
\toprule
\textbf{Etapa} & \textbf{Meta del plan} & \textbf{Avance documentado en el repositorio} \\
\midrule
Semana 1 & Preparar el proyecto, el bucle de juego, Esper, ECS y componentes comunes. & Se establecieron el motor, el mundo Esper, componentes, fábricas, configuración JSON y documentación de contratos. La ventana se configura en 256\,$\times$\,208. \\
Semana 2 & Obtener un prototipo donde se pueda mover y disparar, con cámara y enemigo inicial. & El código incorpora comandos de teclado, movimiento con aceleración, disparos, cámara, enemigos y colisiones. El historial registra la base ECS y la incorporación de la cámara. \\
Semana 3 & Añadir combate, bases, torretas, obstáculos, muerte y puntaje. & Hay sistemas de combate, hitboxes, daño, explosiones, bases, torretas, minas, rocas y vidas del jugador. Los archivos de nivel usan entidades de estas categorías. \\
Semana 4 & Integrar el mundo y la interfaz: cámara, wraparound, estrellas, minimapa, audio y HUD. & El código incorpora cámara, wraparound, campo de estrellas, HUD, minimapa y configuración de música y efectos. En la versión examinada, el minimapa no representa las bases y el HUD mantiene la alerta verde. \\
Semana 5 & Completar diez niveles, escenas y el flujo entre partidas. & El repositorio contiene diez archivos JSON y escenas de título, instrucciones, créditos, pausa, continuación, victoria y derrota. En \archivo{main} aún no se activa el avance de ronda después de destruir todas las bases. \\
Semana 6 & Corregir, publicar, documentar y preparar la presentación. & El repositorio incluye un README y documentos de arquitectura y configuración. En la versión examinada no se encontró un enlace de publicación, una etiqueta final ni el video de presentación. \\
\bottomrule
\end{tabularx}
\end{center}

\section{Hitos verificables en Git}
El historial contiene, entre otros, commits relacionados con la base ECS (\archivo{8212d9c}), compatibilidad con Esper 2 (\archivo{f7ff6a5}), cámara y wraparound (\archivo{6852dc4}), combate (\archivo{2e0eb7d}), controles y módulo del jugador (\archivo{5a2f47a}) y configuración de niveles (\archivo{72fb107}). Estos identificadores sirven para localizar cambios; no bastan por sí solos para atribuir tareas individuales ni para establecer la semana exacta en que se hicieron.

\chapter{Análisis de arquitectura}

\section{Vista general}
\begin{center}
\fbox{\begin{minipage}{0.92\textwidth}
\centering
\textbf{GameEngine / bucle de juego}\par
Procesa eventos, calcula \archivo{delta_time}, actualiza y renderiza\par\vspace{0.6em}
$\downarrow$\par
\textbf{Administrador de escenas} \quad $\longrightarrow$ \quad \textbf{Escena activa}\par
Título, juego, pausa, instrucciones, créditos, continuación, victoria o derrota\par\vspace{0.6em}
$\downarrow$\par
\textbf{Mundo Esper (ECS)}\par
Entidades = identificadores \quad Componentes = datos \quad Sistemas = lógica\par
Fábricas ensamblan jugador, enemigos, bases, torretas, proyectiles y obstáculos\par\vspace{0.6em}
\textbf{Localizador de servicios} \quad $\longrightarrow$ \quad JSON, imágenes, sonidos y fuentes
\end{minipage}}
\end{center}

\section{Bucle principal y escenas}
\archivo{main.py} crea el motor. \archivo{GameEngine.run()} inicializa la escena de título y mantiene un único ciclo que obtiene eventos, calcula el tiempo entre cuadros, actualiza la escena activa, dibuja y presenta cada cuadro. El \archivo{SceneManager} distribuye estas operaciones y aplica los cambios de escena al terminar cada llamada. De este modo, todos los módulos utilizan el mismo ciclo y Pygame permanece activo durante las transiciones.

Las escenas implementan un contrato común para recibir eventos, actualizarse y dibujarse. Esta separación asigna responsabilidades diferenciadas al menú, la partida y las pantallas de resultado. Durante la pausa, la escena de juego deja de ser la escena activa y su actualización queda suspendida.

\section{ECS y responsabilidades}
Esper administra las entidades como identificadores y permite consultarlas mediante sus componentes. Las fábricas de \path{src/create/} construyen entidades; los sistemas de \path{src/ecs/systems/} recorren las combinaciones de componentes que necesitan y actualizan el mundo.

Algunos componentes compartidos son:
\begin{itemize}[leftmargin=*]
  \item \archivo{CTransform}: posición en píxeles de mundo y orientación.
  \item \archivo{CVelocity}: velocidad en píxeles por segundo.
  \item \archivo{CSurface} y \archivo{CAnimation}: representación y animación visual.
  \item \archivo{CHitBox}: caja AABB con tamaño y desplazamiento propios, independiente del sprite.
  \item \archivo{CState}: estado nombrado, tiempo transcurrido y datos opcionales.
  \item Componentes de jugador, combate, bases, torretas, proyectiles y obstáculos.
\end{itemize}

La lógica de movimiento recibe \archivo{delta_time}, por lo que la velocidad se integra en función del tiempo y no solo del número de cuadros. El sistema de combate centraliza detección de solapamientos, daño y destrucción de las entidades relevantes. El mundo envuelve posiciones horizontal y verticalmente mediante el sistema de wraparound.

\section{Patrones y conceptos aplicados}
\begin{description}[leftmargin=!,labelwidth=3.4cm]
  \item[ECS] Separa los datos de las entidades de la lógica de los sistemas. Facilita agregar tipos de entidad combinando componentes y evita concentrar todas las reglas en una jerarquía de clases.
  \item[Command] Los comandos de jugador convierten entradas de teclado en intenciones de movimiento, disparo y habilidad. Las acciones pueden procesarse desde el sistema de comandos sin que cada entidad lea directamente el teclado.
  \item[State] \archivo{CState} representa estados de juego; los sistemas de jugador, enemigos y bases actualizan sus transiciones. En esta implementación los estados se expresan principalmente con datos y lógica de sistemas, no con una clase distinta por estado.
  \item[Service Locator] \archivo{ServiceLocator} ofrece acceso centralizado a configuración, imágenes, sonidos y fuentes, incluyendo caché de recursos.
  \item[Scene] La interfaz \archivo{Scene} y \archivo{SceneManager} organizan pantallas y cambios de flujo bajo el mismo ciclo principal.
  \item[Configuración de datos] JSON define ventana, jugador, enemigos, bases, obstáculos, interfaz, audio y los diez niveles. Esto permite ajustar parámetros sin alterar directamente los sistemas.
\end{description}

\section{Matemáticas, colisiones y respuesta del juego}
El proyecto utiliza vectores bidimensionales para representar dirección, velocidad, cámara y distancia. El movimiento con aceleración y giro, la orientación del sprite y la distancia mínima entre posiciones permiten navegar un espacio con wraparound. Las colisiones usan rectángulos AABB definidos por \archivo{CHitBox}; estas cajas son independientes del tamaño visual del sprite.

La respuesta audiovisual combina sprites animados, disparos, sonidos, explosiones, invulnerabilidad temporal, parpadeo durante la reaparición, interfaz y movimiento de estrellas. Estos elementos comunican al jugador los efectos de sus acciones. La fidelidad del ritmo, la dificultad, la apariencia y el comportamiento debe evaluarse mediante sesiones de comparación con el juego de referencia.

\section{Recursos y configuración}
\archivo{ServiceLocator} localiza y almacena recursos para su reutilización. Los archivos JSON parametrizan el tamaño del mundo, la aparición de enemigos, las bases y los obstáculos. La configuración facilita la ampliación del sistema; sin embargo, cada parámetro requiere un sistema que lo utilice. Por ejemplo, el tiempo de alerta roja está configurado, pero aún no se refleja en un sistema de alertas activo.

\chapter{Evaluación del estado del producto}

La revisión de la rama \archivo{main}, realizada el 23 de noviembre de 2026, permite distinguir las siguientes áreas:
\begin{itemize}[leftmargin=*]
  \item \textbf{Base implementada:} bucle de juego, ECS, escenas, comandos de teclado, cámara, wraparound, jugador, vidas, enemigos, bases, obstáculos y colisiones.
  \item \textbf{Implementación parcial:} transición automática por destrucción de bases; minimapa, pendiente de incorporar marcadores de bases; alertas, cuyo HUD solo muestra el estado verde; puntaje, con récord persistente y vidas extra por completar; y habilidad especial, que debe ajustarse para producir los cuatro proyectiles especificados.
  \item \textbf{Entrega por completar:} verificar el juego completo de principio a fin, publicar la versión, etiquetar el código final y completar documentos y presentación.
\end{itemize}

Esta evaluación documenta el estado del repositorio en la fecha indicada. La validación integral requiere ejecutar una partida completa, revisar los diez niveles y comprobar cada criterio funcional en el juego.

\chapter{Post-mortem grupal}

\section{Aspectos logrados}
La separación por módulos permitió desarrollar en paralelo sobre una arquitectura compartida. Los contratos ECS, las fábricas y la configuración JSON establecieron puntos de integración claros. El historial de Git registra avances incrementales en arquitectura, jugador, combate, recursos y niveles. El uso de \archivo{delta_time}, hitboxes independientes y un administrador común de escenas redujo dependencias entre módulos.

\section{Dificultades y aspectos por mejorar}
La integración requirió coordinar interfaces compartidas: los componentes, las consultas de Esper, las coordenadas, las unidades de velocidad y los contratos de las fábricas deben coincidir entre módulos. En la versión examinada, algunos requisitos configurados aún no se reflejan en el comportamiento del juego; las alertas, el récord, las vidas extra y ciertos marcadores del minimapa son ejemplos de ello. El avance automático entre niveles también debe completarse y verificarse en la rama principal.

\section{Mejoras para proyectos futuros}
Como equipo, proponemos definir desde la primera semana los contratos de componentes, las unidades y las convenciones de coordenadas; integrar módulos pequeños de manera frecuente; y convertir cada criterio obligatorio de la rúbrica en una demostración verificable. También resulta conveniente mantener una matriz de trazabilidad entre requisito, archivo o sistema, responsable y estado. De este modo, una opción de configuración se considerará terminada cuando su efecto esté implementado y verificado en el juego.

\section{Aprendizajes del equipo}
El desarrollo permitió aplicar ECS en un programa interactivo y comprender la importancia de mantener los componentes como datos y los sistemas como lógica. También permitió observar cómo Command y Scene organizan las entradas y las pantallas, mientras que Service Locator centraliza el acceso a recursos. La integración mostró que una arquitectura efectiva requiere contratos claros, validación conjunta y cambios frecuentes que puedan incorporarse al proyecto compartido.

\section{Reflexiones individuales}

\subsection*{Manuel Santiago Prieto Hernández -- Integrante 1}
\textbf{Aporte personal.} Mi responsabilidad se concentró en la arquitectura y la integración del motor: el bucle de juego, la configuración del mundo ECS, los componentes compartidos, la gestión de escenas, los servicios de recursos y la coordinación de la cámara con el mundo. También participé en integrar los módulos para que el jugador, el combate y la interfaz funcionaran dentro del mismo ciclo de actualización.

\textbf{Aspectos logrados.} Fue acertado establecer desde una etapa temprana una estructura común basada en Esper, fábricas y sistemas. El administrador de escenas permitió organizar el menú, la partida, la pausa y las pantallas de resultado sin crear un bucle independiente para cada módulo. La definición de unidades para posición y velocidad, el uso de \archivo{delta_time} y la separación entre hitboxes y sprites proporcionaron contratos concretos para la integración. La cámara y el wraparound permitieron conservar las coordenadas de mundo de las entidades y mantener fija la interfaz.

\textbf{Dificultades.} La integración constituyó el aspecto más exigente de mi responsabilidad. Fue necesario resolver incompatibilidades entre versiones de Esper, diferencias entre las consultas ECS y el desempaquetado de componentes, y rutas de recursos ausentes del repositorio. Estos problemas evidenciaron que un módulo puede funcionar de manera aislada y presentar errores al integrarse con el juego. También confirmé que declarar una opción en JSON no garantiza su funcionamiento: cada parámetro debe conectarse con un sistema y verificarse en ejecución.

\textbf{Oportunidades de mejora.} En un proyecto futuro fijaría las versiones de las dependencias y validaría una instalación limpia desde el inicio. Documentaría los contratos de componentes y las consultas ECS antes de distribuir los sistemas, y realizaría integraciones pequeñas y frecuentes. Asimismo, mantendría una matriz de aceptación que vincule cada requisito con el módulo responsable y con un criterio observable de verificación.

\textbf{Aprendizajes.} Aprendí a modelar el juego mediante datos que los sistemas transforman, en lugar de concentrar todas las reglas en una sola clase. Comprendí cómo el bucle de juego, ECS, Command, State, Scene y Service Locator atienden responsabilidades distintas y complementarias. También consolidé el uso de vectores, tiempo entre cuadros, cámara, coordenadas de mundo y colisiones AABB. Finalmente, comprendí que una arquitectura aporta valor cuando facilita la integración y la revisión del producto completo.

\subsection*{Daniela Munar Gonzalez -- Integrante 2}
\textbf{Aporte y evidencia personal:}\par
\noindent\rule{\textwidth}{0.4pt}\\[1.2em]
\noindent\rule{\textwidth}{0.4pt}\\[1.2em]
\noindent\rule{\textwidth}{0.4pt}\\[1.0em]
\textbf{Qué salió bien:}\par
\noindent\rule{\textwidth}{0.4pt}\\[1.2em]
\noindent\rule{\textwidth}{0.4pt}\\[1.0em]
\textbf{Qué fue difícil y qué cambiaría:}\par
\noindent\rule{\textwidth}{0.4pt}\\[1.2em]
\noindent\rule{\textwidth}{0.4pt}\\[1.0em]
\textbf{Qué aprendí:}\par
\noindent\rule{\textwidth}{0.4pt}\\[1.2em]
\noindent\rule{\textwidth}{0.4pt}

\subsection*{Camilo Puerto Peña -- Integrante 3}
\textbf{Aporte y evidencia personal:}\par
\noindent\rule{\textwidth}{0.4pt}\\[1.2em]
\noindent\rule{\textwidth}{0.4pt}\\[1.2em]
\noindent\rule{\textwidth}{0.4pt}\\[1.0em]
\textbf{Qué salió bien:}\par
\noindent\rule{\textwidth}{0.4pt}\\[1.2em]
\noindent\rule{\textwidth}{0.4pt}\\[1.0em]
\textbf{Qué fue difícil y qué cambiaría:}\par
\noindent\rule{\textwidth}{0.4pt}\\[1.2em]
\noindent\rule{\textwidth}{0.4pt}\\[1.0em]
\textbf{Qué aprendí:}\par
\noindent\rule{\textwidth}{0.4pt}\\[1.2em]
\noindent\rule{\textwidth}{0.4pt}

\subsection*{Camilo Inoue Ruda -- Integrante 4}
\textbf{Aporte y evidencia personal:}\par
\noindent\rule{\textwidth}{0.4pt}\\[1.2em]
\noindent\rule{\textwidth}{0.4pt}\\[1.2em]
\noindent\rule{\textwidth}{0.4pt}\\[1.0em]
\textbf{Qué salió bien:}\par
\noindent\rule{\textwidth}{0.4pt}\\[1.2em]
\noindent\rule{\textwidth}{0.4pt}\\[1.0em]
\textbf{Qué fue difícil y qué cambiaría:}\par
\noindent\rule{\textwidth}{0.4pt}\\[1.2em]
\noindent\rule{\textwidth}{0.4pt}\\[1.0em]
\textbf{Qué aprendí:}\par
\noindent\rule{\textwidth}{0.4pt}\\[1.2em]
\noindent\rule{\textwidth}{0.4pt}

\chapter{Roles y contribuciones del equipo}
\small
\setlength{\tabcolsep}{4pt}
\begin{tabular}{@{}>{\raggedright\arraybackslash}p{0.13\textwidth} >{\raggedright\arraybackslash}p{0.26\textwidth} >{\raggedright\arraybackslash}p{0.50\textwidth}@{}}
\toprule
\textbf{Integrante} & \textbf{Responsabilidad acordada} & \textbf{Contribución que se debe describir con evidencia propia} \\
\midrule
1 & Manuel Santiago Prieto Hernández -- arquitectura, motor y ECS. & Motor, loop, escenas, componentes compartidos, Service Locator, cámara e integración. \\
2 & Daniela Munar Gonzalez -- jugador. & Controles, movimiento, disparos, habilidad, vidas y reaparición. \\
3 & Camilo Puerto Peña -- combate. & Enemigos, bases, torretas, obstáculos, hitboxes, colisiones y explosiones. \\
4 & Camilo Inoue Ruda -- mundo e interfaz. & Estrellas, minimapa, HUD, alertas, puntaje, niveles, audio y presentación. \\
\bottomrule
\end{tabular}
\normalsize
\setlength{\tabcolsep}{6pt}

\chapter{Conclusiones}
La arquitectura ofrece una base coherente para una recreación de Bosconian: un solo loop, entidades compuestas mediante componentes, sistemas especializados, escenas, recursos centralizados y niveles configurables. El trabajo pendiente más importante consiste en cerrar las brechas funcionales identificadas y validar el recorrido completo con el juego en ejecución. La entrega debe acompañarse de una versión reproducible, el código final etiquetado, la publicación y las evidencias de participación de cada integrante.

\chapter*{Fuentes de trabajo}
\addcontentsline{toc}{chapter}{Fuentes de trabajo}
\begin{itemize}[leftmargin=*]
  \item Enunciado y rúbrica del proyecto suministrados por el curso.
  \item Repositorio del equipo: \url{https://github.com/grupo-amistad/bosconian-recreation}.
  \item Documentación de arquitectura: \path{docs/arquitectura.md}.
  \item Plan de trabajo: \path{docs/plan-de-trabajo.md}.
  \item Esquema de niveles: \path{docs/configuracion-de-niveles.md}.
  \item Código del motor, ECS, jugador, servicios e interfaz en sus directorios dentro de \path{src/}.
  \item Configuración de ventana, entidades, audio y niveles: archivos JSON en \path{assets/cfg/}.
\end{itemize}

\end{document}
